\subsection{The Wonderful Exchange}\label{sec:exchange}

On the octave day of Christmas the Roman Breviary began Lauds, the Hour that is
now Morning Prayer, with an antiphon that names the whole mystery in three words,
\latin{O admirabile commercium}, ``How wonderful is the dealing here!'' It goes
on:

\begin{quote}
The Maker of man hath taken to Himself a living body, and hath been pleased to
be born of the Virgin: He is come forth conceived without seed, and hath made us
partakers of His Divine nature.\par
{\smallskip\noindent\small ---Antiphon \latin{O admirabile commercium}, the Roman Breviary, in Bute's English\par}
\end{quote}

\noindent An exchange is a trade: each party gives what is his and receives what
is the other's. At Bethlehem God takes what is ours, a body, a birth, a mother,
and gives what is his. The Fathers never tired of the trade. On the 18th of
December the Office reads the \emph{Letter to Diognetus}, and its unknown
author breaks off his argument to cry out:

\begin{quote}
For what other thing was capable of covering our sins than His righteousness?
By what other one was it possible that we, the wicked and ungodly, could be
justified, than by the only Son of God? O sweet exchange! O unsearchable
operation! O benefits surpassing all expectation! that the wickedness of many
should be hid in a single righteous One, and that the righteousness of One
should justify many transgressors!\par
{\smallskip\noindent\small ---Letter to Diognetus, 9\par}
\end{quote}

\noindent St Gregory Nazianzen, read on the Tuesday of the first week, set the
two sides of the trade against each other:

\begin{quote}
And He who gives riches becomes poor; for He assumes the poverty of my flesh,
that I may assume the riches of His Godhead. He that is full empties Himself;
for He empties Himself of His Glory for a short while, that I may have a share
in His Fulness. \ldots{} I had a share in the Image and I did not keep it; He
partakes of my flesh that He may both save the Image and make the flesh
immortal. He communicates a Second Communion, far more marvellous than the
first, inasmuch as then He imparted the better nature, but now He Himself
assumes the worse.\par
{\smallskip\noindent\small ---St Gregory Nazianzen, Oration 45\par}
\end{quote}

\noindent St Irenaeus had put it in a single line: the
Word of God, our Lord Jesus Christ, ``did, through His transcendent love, become
what we are, that He might bring us to be even what He is Himself.'' St
Athanasius wrote of the Word made flesh: ``For He was made man that we might be
made God.'' St Ambrose, as the Catena Aurea gives him, turned the exchange at
the manger into a prayer of thanks: ``On thy account then am I weak, in thee am I strong. On thy account
am I poor, in thee am I rich.''

Pope St Leo the Great, preaching on the Passion, carried the trade to the
Cross: ``For He had come into this world a rich and merciful Merchant from the
skies, and by a wondrous exchange had entered into a bargain of salvation with
us, receiving ours and giving His, honour for insults, salvation for pain, life
for death.''

The Church sings the exchange at the altar. The third Preface of the Nativity is
a reworking of a December preface from the Veronese Sacramentary, which gave
thanks because \latin{magnificum nostrae commercium reparationis effulsit},
``the magnificent exchange of our restoration has shone forth,'' and ended with the
gift received in return: \latin{nos quoque mirando consortio reddis aeternos},
``thou makest us also eternal by a wondrous fellowship.'' Through the Christmas
season the priest may give thanks in this Preface, and the Church on earth takes
up the old trade in her Eucharist: she offers bread and wine, and receives the
Body and Blood of the Lord.

St John of the Cross, in his ballad of the Nativity, set the exchange at the
manger, and gave the wonder to the Mother:

\begin{verse}
So the bride at her betrothal\\
Did the bridal gifts arrange;\\
But the Mother looked in wonder\\
At the marvellous exchange.\\
Man gave forth a song of gladness,\\
God Himself a plaintive moan;\\
Both possessing that which never\\
Had been hitherto their own.
\end{verse}
